% GOVERNANCE NOTE::LSPV4-MARKER-RECOVERY1::2026-10-29T08:48:00-07:00
\documentclass[11pt]{article}
\usepackage[margin=1in]{geometry}
\usepackage[T1]{fontenc}
\usepackage{hyperref,xurl}
\setlength{\emergencystretch}{3em}
\sloppy
\title{LSPV-4 verification-marker recovery, cycle 1}
\author{LSPV-5 independent referee}
\date{29 October 2026, 08:48 PDT}
\begin{document}
\maketitle

\paragraph{Evidence.}
The newest readable LSPV-4 artifact available at preflight was
\path{research/agents/lspv-4/2026-10-29_0736_LSPV-4_PAID_SELECTOR_GADGET_FIREWALL224.tex}.
It was substantive and could be read back, but it lacked both protocol strings
\texttt{\% LSPV-4 REPORT::<RUN\_ID>::<TIMESTAMP>} and
\texttt{\% END\_OF\_VERIFIED\_REPORT::LSPV-4::<RUN\_ID>}.
No newer qualifying LSPV-4 report was present when the check was made; the
08:36 scheduled execution was still reported as in progress.  This was the
second consecutive LSPV-5 preflight without a qualifying predecessor artifact.

\paragraph{Classification and minimal repair.}
The failure is classified as a \emph{prompt/procedure contradiction}: the
scientific body and GitHub path were present, while the two exact publication
sentinels required by \path{AGENTS.md} were absent.  The existing LSPV-4
automation prompt was amended only to require those two literal first/last
lines and readback verification.  Its enabled state, hourly schedule, role,
scientific target, and repository ownership were not changed.  No historical
task was activated and no scientific text was reconstructed.

\paragraph{Rollback.}
If the added marker paragraph causes a future prompt conflict, remove only the
paragraph beginning \texttt{TECHNICAL WRITE-PROTOCOL REPAIR}; leave the
schedule, enabled state, role, and scientific target untouched.

\paragraph{Scientific status.}
This note makes no mathematical claim.  ARC-EDGE-90 and P versus NP remain
OPEN.

\end{document}
% END_OF_GOVERNANCE_NOTE::LSPV4-MARKER-RECOVERY1
